% !TeX root = ../notes.tex
% Planning version: section skeleton, content bullets, live TikZ, figure list, and reading list.

\chapter{From automated synthesis to self-driving chemistry}

\begin{notebox}
  \keypoint{Intellectual question:} How is the loop closed between experiment selection, execution, measurement, and learning?
\end{notebox}

\section{Automation, programmability, and autonomy}

\begin{itemize}
  \item Define automation as execution of known instructions.
  \item Define programmability as a digital representation that can be transferred, generalised, and executed on compatible hardware.
  \item Define autonomy as measurement changing the decision about what happens next.
  \item Keep these categories separate when evaluating published platforms.
\end{itemize}

\solution{5.8cm}{%
  \centering
  \resizebox{0.94\linewidth}{!}{%
    \begin{tikzpicture}[
      stage/.style={draw=darkBlue,fill=paleBlue,line width=1.6pt,minimum width=4.0cm,minimum height=1.2cm,align=center,font=\sffamily},
      final/.style={draw=darkGreen,fill=paleGreen,line width=1.8pt,minimum width=4.0cm,minimum height=1.2cm,align=center,font=\sffamily},
      arrow/.style={->,>=Stealth,draw=darkGrey,line width=2.2pt}
    ]
      \node[stage] (automation) at (0,0) {automation\\executes instructions};
      \node[stage] (program) at (5.5,0) {programmability\\encodes procedures};
      \node[final] (autonomy) at (11,0) {autonomy\\uses results to decide};
      \draw[arrow] (automation) -- (program);
      \draw[arrow] (program) -- (autonomy);
    \end{tikzpicture}%
  }
}

\section{Automated peptide synthesis}

\begin{itemize}
  \item Use Merrifield and Stewart as the historical example of mechanising a repetitive synthetic cycle.\mcite{Merrifield1965}
  \item Identify the predefined sequence of deprotection, washing, coupling, and washing.
  \item Credit the scientific importance of reliable execution without calling it autonomous.
  \item Ask what is represented digitally and which choices remain fixed before the run begins.
\end{itemize}

\section{Chemputer and programmable synthesis}

\begin{itemize}
  \item Introduce transferable digital procedures and a hardware graph.\mcite{Steiner2019}
  \item Show how reaction, work-up, isolation, and purification become machine-executable operations.
  \item Separate procedural abstraction from adaptive experiment selection.
  \item Use diphenhydramine, rufinamide, and sildenafil only as evidence that the language controls multistep synthesis.
\end{itemize}

\solution{6.5cm}{%
  \centering
  \resizebox{0.94\linewidth}{!}{%
    \begin{tikzpicture}[
      box/.style={draw=darkBlue,fill=paleBlue,line width=1.5pt,minimum width=3.1cm,minimum height=1.0cm,align=center,font=\sffamily},
      arrow/.style={->,>=Stealth,draw=darkGrey,line width=2pt},
      label/.style={font=\small\sffamily,text=darkGrey,align=center}
    ]
      \node[box] (procedure) at (0,0) {written chemical\\procedure};
      \node[box] (code) at (4.4,0) {chemical\\program};
      \node[box] (graph) at (8.8,0) {hardware graph\\and compiler};
      \node[box] (execute) at (13.2,0) {physical unit\\operations};
      \draw[arrow] (procedure) -- (code);
      \draw[arrow] (code) -- (graph);
      \draw[arrow] (graph) -- (execute);
      \node[label] at (6.6,-1.2) {programmability makes procedures transferable};
    \end{tikzpicture}%
  }
}

\section{From automated experimentation to a closed loop}

\begin{itemize}
  \item Identify sensors, analytical instruments, orchestration, decision logic, and model updating.
  \item State that the loop is only closed when a measurement changes a subsequent experimental action.
  \item Include data-quality gates, instrument status, and safe fallback behaviour.
  \item Reuse Chapter~6's experiment-selection ideas and Chapter~7's execution and measurement infrastructure.
\end{itemize}

\solution{8.0cm}{%
  \centering
  \resizebox{0.86\linewidth}{!}{%
    \begin{tikzpicture}[
      stage/.style={draw=darkBlue,fill=paleBlue,line width=1.6pt,minimum width=3.0cm,minimum height=1.0cm,align=center,font=\sffamily},
      update/.style={draw=darkGreen,fill=paleGreen,line width=1.8pt,minimum width=3.0cm,minimum height=1.0cm,align=center,font=\sffamily},
      arrow/.style={->,>=Stealth,draw=darkGrey,line width=2.2pt}
    ]
      \node[stage] (select) at (0,3.0) {select};
      \node[stage] (execute) at (4.4,3.0) {execute};
      \node[stage] (measure) at (4.4,0) {measure and QC};
      \node[update] (update) at (0,0) {update knowledge};
      \draw[arrow] (select) -- (execute);
      \draw[arrow] (execute) -- (measure);
      \draw[arrow] (measure) -- (update);
      \draw[arrow] (update) -- (select);
    \end{tikzpicture}%
  }
}

\section{Cooper: a mobile robotic chemist}

\begin{itemize}
  \item Use the mobile robot as an orchestration layer between existing laboratory instruments.\mcite{Burger2020}
  \item Show the ten-variable photocatalyst-formulation search and batched Bayesian selection.
  \item Discuss 688 experiments over eight days and the improvement over initial formulations.
  \item Distinguish mobility, automation, and autonomous experiment selection.
\end{itemize}

\section{RoboChem}

\begin{itemize}
  \item Use the platform as the clearest true closed-loop optimisation case.\mcite{Slattery2024}
  \item Identify flow photochemistry, phase sensing, inline NMR, Bayesian optimisation, and multiobjective choice.
  \item Show optimisation, intensification, and scale-up as linked but distinct goals.
  \item Ask which objectives and constraints are encoded and which chemical judgements remain external.
\end{itemize}

\section{Autonomous exploratory synthesis}

\begin{itemize}
  \item Use the 2024 Nature platform as the final course case.\mcite{Dai2024}
  \item Show mobile robots connecting synthesis, UPLC-MS, and benchtop NMR.
  \item Explain the heuristic planner, orthogonal analytical thresholds, replication, scale-up, diversification, and functional assays.
  \item State explicitly that the demonstrated workflows are autonomous but are not closed-loop optimisation processes.
  \item Use the unexpected cyclisation missed by the initial UPLC-MS decision rule to expose analytical and contextual limits.
\end{itemize}

\solution{7.2cm}{%
  \centering
  \resizebox{0.92\linewidth}{!}{%
    \begin{tikzpicture}[
      scalar/.style={draw=darkBlue,fill=paleBlue,line width=1.5pt,minimum width=4.4cm,minimum height=1.2cm,align=center,font=\sffamily},
      explore/.style={draw=darkGreen,fill=paleGreen,line width=1.7pt,minimum width=4.4cm,minimum height=1.2cm,align=center,font=\sffamily},
      arrow/.style={->,>=Stealth,draw=darkGrey,line width=2pt},
      label/.style={font=\small\sffamily,text=darkGrey,align=center}
    ]
      \node[scalar] (opt) at (0,1.4) {scalar optimisation\\improve a defined objective};
      \node[explore] (search) at (0,-1.4) {exploratory autonomy\\screen, classify, replicate, branch};
      \node[scalar] (loop) at (7.0,1.4) {model update and\\next-point selection};
      \node[explore] (rules) at (7.0,-1.4) {heuristic rules and\\orthogonal thresholds};
      \draw[arrow] (opt) -- (loop);
      \draw[arrow] (search) -- (rules);
      \node[label] at (3.5,-2.4) {both are autonomous, but they answer different chemical questions};
    \end{tikzpicture}%
  }
}

\section{Robustness, error handling, and recovery}

\begin{itemize}
  \item Distinguish hardware faults, analytical anomalies, and chemically valid negative outcomes.
  \item Include calibration checks, watchdogs, inventory tracking, reruns, and safe shutdown.
  \item Require a record of why an experiment was rejected, repeated, or escalated to a human.
  \item Show that autonomy without trustworthy failure handling simply accelerates unrecognised error.
\end{itemize}

\solution{7.0cm}{%
  \centering
  \resizebox{0.9\linewidth}{!}{%
    \begin{tikzpicture}[
      decision/.style={draw=darkBlue,fill=paleBlue,line width=1.5pt,minimum width=3.4cm,minimum height=1.0cm,align=center,font=\sffamily},
      pass/.style={draw=darkGreen,fill=paleGreen,line width=1.5pt,minimum width=3.4cm,minimum height=1.0cm,align=center,font=\sffamily},
      fail/.style={draw=darkRed,fill=paleRed,line width=1.5pt,minimum width=3.4cm,minimum height=1.0cm,align=center,font=\sffamily},
      arrow/.style={->,>=Stealth,draw=darkGrey,line width=2pt}
    ]
      \node[decision] (qc) at (0,0) {measurement passes QC?};
      \node[pass] (update) at (6,1.5) {update and continue};
      \node[fail] (diagnose) at (6,-1.5) {diagnose: instrument,\\analysis, or chemistry};
      \node[fail] (human) at (12,-1.5) {retry, safe stop,\\or human escalation};
      \draw[arrow] (qc) -- node[above]{yes} (update);
      \draw[arrow] (qc) -- node[below]{no} (diagnose);
      \draw[arrow] (diagnose) -- (human);
    \end{tikzpicture}%
  }
}

\section{Human-in-the-loop and responsible autonomy}

\begin{itemize}
  \item Humans define the chemical problem, objectives, constraints, priors, and success criteria.
  \item The system may select experiments only within that delegated decision space.
  \item Discuss explainability at the level needed to audit why an experiment was chosen.
  \item Include safety, resource use, reproducibility, and the cost of analysis in objective design.
  \item Do not imply that automation removes the need for chemical judgement.
\end{itemize}

\section{Closing the course}

\begin{itemize}
  \item Reconnect the full progression: describe, model, choose, execute, autonomise.
  \item Close with $\text{select}\rightarrow\text{execute}\rightarrow\text{measure}\rightarrow\text{update}\rightarrow\text{select}$.
  \item Emphasise that every loop inherits the limitations of representation, data, measurement, model, and objective.
  \item End with digital chemistry as scalable chemical reasoning rather than replacement of the chemist.
\end{itemize}

\begin{insightbox}
  \keypoint{A self-driving laboratory is not one instrument.}
  It is a tested relationship between a chemical objective, executable operations, trustworthy measurements, explicit decision logic, and safe recovery.
\end{insightbox}

\section{Summary and review}

\begin{itemize}
  \item Which operations are automated?
  \item Is the procedure digitally programmable and transferable?
  \item Which measurement changes the next decision?
  \item Is the platform optimising a scalar objective or exploring several possible outcomes?
  \item What happens when hardware, analysis, or chemistry fails?
\end{itemize}

\section*{Author production checklist: publication figures}

\begin{itemize}
  \item Merrifield and Stewart: the authentic synthesiser/process figure to establish repetitive predefined automation.
  \item Steiner \textit{et al.}: Fig.~1 for abstraction and Chemputer architecture; Fig.~2 for information flow from written procedure and hardware graph to commands.
  \item Burger \textit{et al.}: the platform/workflow figure and the search-trajectory/result figure from the final Nature PDF.
  \item Slattery \textit{et al.}: Fig.~1 for the complete RoboChem concept; Fig.~2 for platform architecture and phase tracking; Fig.~3 for a representative optimisation campaign.
  \item Dai \textit{et al.}: Fig.~1 for the modular laboratory and heuristic planner; Fig.~2 for divergent synthesis; Fig.~3 for autonomous host--guest discovery; Fig.~4 only if the photoreaction extension is discussed.
\end{itemize}

\section*{Author reading and PDF-download list}

\begin{itemize}
  \item Historical automated peptide synthesis.\mcite{Merrifield1965}
  \item Chemputer and chemical programming language.\mcite{Steiner2019}
  \item Mobile robotic chemist.\mcite{Burger2020}
  \item RoboChem closed-loop flow photochemistry.\mcite{Slattery2024}
  \item Autonomous exploratory synthetic chemistry.\mcite{Dai2024}
\end{itemize}
